\scanpage{56}
\margin{Ceci donne une explication}

\uline{Conj.} Si $X$ loc. de t.f. sur $\mathbb{C}$, alors
\[
H^*_{\substack{\unsure{\mathrm{Str}}\\\text{ou crist.}}}
(X/\mathbb{C},\underline{\mathcal{O}}_X)
\simeq H^*(X^{\mathrm{an}},\mathbb{C})
\]

N.B. La conjecture est prouvée, grâce au th, \deleted{lorsque}, et au th. de comparaison de l'exposé 1, dans le cas où $X$ est lisse sur $\mathbb{C}$.
Donc \uline{dans le cas lisse, au dessus d'une base de} car. 0, \uline{on trouve une description de la coh. de} D.R. \uline{sans utiliser de différentielles}.

\margin{55}
